% Guide de mise en ligne — Fils Naturel Tech-Com
% Compilation : pdflatex hebergement-fntc.tex (deux passes pour la table des matières)

\documentclass[11pt,a4paper]{article}

\usepackage{lmodern}
\usepackage[T1]{fontenc}
\usepackage[utf8]{inputenc}
\usepackage[french]{babel}
\usepackage[a4paper,margin=2.4cm,headheight=15pt]{geometry}
\usepackage{xcolor}
\usepackage{titlesec}
\usepackage{enumitem}
\usepackage{fancyhdr}
\usepackage{tabularx}
\usepackage{booktabs}
\usepackage{longtable}
\usepackage{tcolorbox}
\usepackage[hidelinks]{hyperref}

\definecolor{sig}{HTML}{D8161F}
\definecolor{amb}{HTML}{E8B33A}
\definecolor{encre}{HTML}{1A0D0C}
\definecolor{encre2}{HTML}{6B5D57}
\definecolor{os}{HTML}{F4F2EC}
\definecolor{os2}{HTML}{EAE7DE}
\definecolor{vert}{HTML}{3F8C5E}

\hypersetup{
  pdftitle={Guide de mise en ligne — Fils Naturel Tech-Com},
  pdfauthor={Fils Naturel Tech-Com},
  pdfsubject={Acheter un domaine, payer l'hébergement, mettre le site en ligne},
}

\pagestyle{fancy}
\fancyhf{}
\fancyhead[L]{\small\color{encre2}Fils Naturel Tech-Com}
\fancyhead[R]{\small\color{encre2}Guide de mise en ligne}
\fancyfoot[C]{\small\color{encre2}\thepage}
\renewcommand{\headrulewidth}{0.4pt}
\renewcommand{\headrule}{\hbox to\headwidth{\color{os2}\leaders\hrule height \headrulewidth\hfill}}

\titleformat{\section}{\Large\bfseries\color{encre}}{\thesection.}{0.6em}{}
\titleformat{\subsection}{\large\bfseries\color{encre}}{\thesubsection}{0.6em}{}
\titlespacing*{\section}{0pt}{1.6em}{0.7em}

\setlist[itemize]{leftmargin=1.2em,itemsep=0.25em,topsep=0.35em}
\setlist[enumerate]{leftmargin=1.5em,itemsep=0.4em,topsep=0.4em}

\newtcolorbox{conseil}[1][]{colback=os,colframe=amb,boxrule=0pt,leftrule=3pt,
  arc=2pt,left=8pt,right=8pt,top=6pt,bottom=6pt,#1}
\newtcolorbox{attention}[1][]{colback=os,colframe=sig,boxrule=0pt,leftrule=3pt,
  arc=2pt,left=8pt,right=8pt,top=6pt,bottom=6pt,#1}
\newtcolorbox{etape}[1][]{colback=os2,colframe=encre,boxrule=0pt,leftrule=3pt,
  arc=2pt,left=8pt,right=8pt,top=6pt,bottom=6pt,#1}

\newcommand{\bouton}[1]{\textbf{\textcolor{sig}{[\,#1\,]}}}
\newcommand{\champ}[1]{\textit{#1}}
\newcommand{\saisir}[1]{\texttt{#1}}
\newcommand{\lien}[2]{\href{#1}{\texttt{#2}}}

% Ligne à remplir à la main.
\newcommand{\aremplir}{\dotfill\rule{0pt}{1.1em}}

\begin{document}

\begin{titlepage}
\centering
\vspace*{2.5cm}
{\color{sig}\rule{4cm}{2pt}}\par
\vspace{1.2cm}
{\Huge\bfseries Mettre le site en ligne\par}
\vspace{0.9cm}
{\LARGE Fils Naturel Tech-Com\par}
\vspace{2.2cm}
{\large De l'achat du nom de domaine\\
jusqu'au site ouvert au public.\par}
\vspace{1cm}
{\color{encre2}\large Écrit pour quelqu'un qui n'a jamais fait ça.\\
Aucune connaissance en informatique n'est nécessaire.\par}
\vfill
{\color{encre2}\large Version 1 — \today\par}
\vspace{0.5cm}
{\color{sig}\rule{4cm}{2pt}}\par
\end{titlepage}

\tableofcontents
\newpage

% ==================================================================
\section{À lire avant tout}

\subsection{Ce que vous allez faire}

Le site est écrit, testé, et son code est rangé sur Internet. Il lui manque
trois choses, que ce guide vous fait acheter et brancher :

\begin{enumerate}
  \item \textbf{Un nom de domaine} — l'adresse que les clients taperont, par
        exemple \saisir{filsnaturel.com}. On l'achète une fois, on le
        renouvelle chaque année.
  \item \textbf{Un hébergement} — l'ordinateur qui fait tourner le site jour et
        nuit. On le paie chaque mois.
  \item \textbf{Une base de données} — l'endroit où sont rangés les articles,
        les commandes et les clients. On la paie chaque mois aussi.
\end{enumerate}

\begin{conseil}
\textbf{Une image simple.} Le nom de domaine, c'est l'adresse de la boutique
écrite sur la rue. L'hébergement, c'est le local. La base de données, c'est le
stock et les registres à l'intérieur. Sans local, l'adresse ne mène nulle part.
\end{conseil}

\subsection{Combien de temps, combien d'argent}

Comptez \textbf{une heure} en tout, sans compter l'attente automatique du
branchement du domaine (de dix minutes à quelques heures, pendant lesquelles
vous n'avez rien à faire).

\begin{center}
\begin{tabularx}{\textwidth}{@{}Xrr@{}}
\toprule
\textbf{Ce que vous payez} & \textbf{Par mois} & \textbf{Par an} \\
\midrule
Nom de domaine \saisir{.com} & --- & 10 à 15 \$ \\
Hébergement du site (Render, offre \champ{Starter}) & 7 \$ & 84 \$ \\
Base de données (Render PostgreSQL, offre \champ{Basic}) & 6 \$ & 72 \$ \\
Stockage de la base (au-delà du forfait) & $\approx$ 0,30 \$ & $\approx$ 4 \$ \\
\midrule
\textbf{Total} & \textbf{$\approx$ 13 \$} & \textbf{$\approx$ 175 \$} \\
\bottomrule
\end{tabularx}
\end{center}

Environ \textbf{175 dollars la première année}, puis autant les années
suivantes. Les prix sont ceux constatés en 2026 ; vérifiez-les sur
\lien{https://render.com/pricing}{render.com/pricing} avant de payer.

\begin{attention}
\textbf{Ne restez pas sur l'offre gratuite de Render.} La base de données
gratuite est \textbf{effacée au bout de 30 jours}. Tout disparaît : articles,
photos, commandes, clients, comptes. L'offre gratuite sert à essayer, pas à
tenir une boutique.
\end{attention}

\subsection{Ce qu'il faut avoir sous la main}

\begin{itemize}
  \item Une adresse email à vous, que vous consultez vraiment.
  \item Une \textbf{carte bancaire Visa ou Mastercard} acceptée à
        l'international, ou un compte PayPal alimenté.
  \item Un ordinateur (pas un téléphone : certaines pages sont trop denses).
  \item Une heure au calme.
\end{itemize}

\begin{conseil}
\textbf{Si vous n'avez pas de carte internationale.} Deux solutions courantes :
une carte prépayée virtuelle (proposée par plusieurs banques et applications de
transfert), ou demander à un proche disposant d'une carte de payer pour vous.
Dans ce dernier cas, \textbf{créez quand même les comptes à votre nom et avec
votre email} : seule la carte est empruntée. Les comptes doivent rester les
vôtres.
\end{conseil}

\subsection{Les mots que vous allez croiser}

\begin{center}
\begin{tabularx}{\textwidth}{@{}lX@{}}
\toprule
\textbf{Mot} & \textbf{Ce que ça veut dire} \\
\midrule
Domaine & L'adresse du site : \saisir{filsnaturel.com}. \\
Registrar & L'entreprise chez qui on achète le domaine. \\
DNS & L'annuaire qui relie l'adresse à l'ordinateur. C'est ce qu'on règle à l'étape~6. \\
Enregistrement A & Une ligne de l'annuaire qui donne le numéro (l'IP) de l'ordinateur. \\
Enregistrement CNAME & Une ligne qui dit « cette adresse est un autre nom pour celle-là ». \\
HTTPS / SSL & Le cadenas dans le navigateur. Render s'en occupe, gratuitement. \\
Déploiement & L'opération qui met la dernière version du site en ligne. \\
Variable d'environnement & Un réglage secret donné au site, comme un mot de passe. \\
\bottomrule
\end{tabularx}
\end{center}

% ==================================================================
\section{Étape 1 — Choisir le nom de domaine}

\subsection{Comment choisir}

Le nom sera dit au téléphone, écrit sur une carte, tapé de mémoire. Donc :

\begin{itemize}
  \item \textbf{Court}. \saisir{filsnaturel.com} vaut mieux que \saisir{filsnatureltechcomhaiti.com}.
  \item \textbf{Sans tiret ni chiffre}. « Tiret du six » au téléphone, c'est perdu d'avance.
  \item \textbf{Facile à épeler} pour quelqu'un qui l'entend pour la première fois.
  \item \textbf{En \saisir{.com}} de préférence : c'est ce que les gens tapent par réflexe.
\end{itemize}

Quelques pistes à tester :

\begin{itemize}
  \item \saisir{filsnaturel.com}
  \item \saisir{fntc.com}
  \item \saisir{fntchaiti.com}
  \item \saisir{filsnaturelhaiti.com}
\end{itemize}

\begin{conseil}
Le \saisir{.ht} est le domaine d'Haïti. Il est plus cher, plus long à obtenir,
et beaucoup de gens tapent \saisir{.com} par habitude. Prenez un \saisir{.com}.
\end{conseil}

\subsection{Vérifier qu'il est libre}

Allez sur \lien{https://www.namecheap.com}{namecheap.com}, tapez le nom voulu
dans la barre de recherche, appuyez sur Entrée. Le site répond \champ{available}
(libre) ou \champ{taken} (déjà pris). S'il est pris, essayez une variante.

\begin{attention}
Ne payez pas les extras proposés à la volée : \champ{hosting},
\champ{email hosting}, \champ{SSL certificate}, \champ{Premium DNS}. Vous n'en
avez besoin d'aucun. L'hébergement, c'est Render ; le certificat SSL, Render le
donne gratuitement.
\end{attention}

% ==================================================================
\section{Étape 2 — Acheter le domaine}

\subsection{Où l'acheter}

\begin{center}
\begin{tabularx}{\textwidth}{@{}lXr@{}}
\toprule
\textbf{Vendeur} & \textbf{Pourquoi} & \textbf{Prix \saisir{.com}} \\
\midrule
\lien{https://www.namecheap.com}{Namecheap} & Le plus simple pour un débutant. Accepte PayPal. Vie privée incluse. & 12 à 15 \$/an \\
\lien{https://www.cloudflare.com/products/registrar/}{Cloudflare} & Le moins cher (prix coûtant, aucune marge). Interface plus austère. & 10 à 11 \$/an \\
\lien{https://www.ovhcloud.com}{OVH} & Support en français, facturation en euros. & 12 à 15 \$/an \\
\bottomrule
\end{tabularx}
\end{center}

\textbf{Recommandation : Namecheap}, pour une première fois. La suite du guide
le prend en exemple ; les autres fonctionnent pareil, avec d'autres mots.

\begin{attention}
Méfiez-vous des offres « \saisir{0,99 \$} la première année ». Le renouvellement
grimpe souvent à 20 ou 30 \$. Regardez toujours le prix de \textbf{renouvellement},
pas celui de la première année.
\end{attention}

\subsection{Acheter, pas à pas}

\begin{etape}
\textbf{Sur \lien{https://www.namecheap.com}{namecheap.com}}
\end{etape}

\begin{enumerate}
  \item En haut à droite, \bouton{Sign Up} : créez un compte avec votre email.
        Notez bien le mot de passe, vous en aurez besoin chaque année.
  \item Cherchez votre nom de domaine, puis \bouton{Add to Cart}.
  \item Cliquez sur \bouton{Checkout}.
  \item Sur la page du panier, réglez :
  \begin{itemize}
    \item \champ{Domain Registration} : \textbf{1 an} suffit (ou 2 ans si vous
          préférez ne pas y penser).
    \item \champ{Auto-Renew} : \textbf{activé}. Sans cela, le domaine expire et
          quelqu'un d'autre peut le prendre.
    \item \champ{Domain Privacy} (WithheldForPrivacy) : \textbf{activé}. C'est
          gratuit chez Namecheap et ça cache votre adresse personnelle des
          annuaires publics.
    \item Tout le reste : \textbf{décoché}.
  \end{itemize}
  \item \bouton{Confirm Order}, puis payez par carte ou PayPal.
  \item Ouvrez l'email de confirmation et \textbf{cliquez sur le lien de
        vérification}. Sans cette validation, le domaine est suspendu au bout de
        quelques jours.
\end{enumerate}

\begin{conseil}
\textbf{Notez tout de suite} : l'adresse email du compte Namecheap, le mot de
passe, le nom de domaine exact, la date d'achat. La fiche du
chapitre~\ref{sec:fiche} est faite pour ça.
\end{conseil}

% ==================================================================
\section{Étape 3 — Créer le compte d'hébergement}

\begin{etape}
\textbf{Sur \lien{https://render.com}{render.com}}
\end{etape}

\begin{enumerate}
  \item \bouton{Get Started} ou \bouton{Sign Up}.
  \item Choisissez \bouton{Sign up with GitHub} : c'est là qu'est rangé le code
        du site, et cette liaison évite bien des manipulations.
  \item Si vous n'avez pas de compte GitHub, créez-le d'abord sur
        \lien{https://github.com}{github.com} — c'est gratuit — puis revenez.
  \item GitHub demande d'autoriser Render à lire vos dépôts : acceptez.
  \item Confirmez votre adresse email en cliquant le lien reçu.
\end{enumerate}

\begin{attention}
Le compte Render doit être créé \textbf{à votre nom et avec votre email}. C'est
lui qui détient le site et la base. Ne le laissez pas au nom d'un tiers.
\end{attention}

% ==================================================================
\section{Étape 4 — Créer la base de données}

On crée la base \textbf{avant} le site : le site a besoin de son adresse pour
démarrer.

\begin{etape}
\textbf{Dans le tableau de bord Render}
\end{etape}

\begin{enumerate}
  \item \bouton{New +} en haut à droite, puis \bouton{Postgres}.
  \item \champ{Name} : \saisir{fntc-base} (ou ce que vous voulez).
  \item \champ{Region} : choisissez \textbf{Ohio} ou \textbf{Virginia} — les
        plus proches d'Haïti, donc les plus rapides.
  \item \champ{PostgreSQL Version} : laissez la valeur proposée.
  \item \champ{Instance Type} : choisissez \textbf{Basic-256mb}, à environ
        \textbf{6 \$ par mois}. \textbf{Surtout pas \champ{Free}} : elle serait
        effacée au bout de 30 jours.
  \item \bouton{Create Database}. Render demande une carte bancaire si ce n'est
        pas déjà fait : saisissez-la, elle servira aussi pour le site.
  \item Attendez que le statut passe à \champ{Available} (une ou deux minutes).
\end{enumerate}

\subsection{Récupérer l'adresse de la base}

Sur la page de la base, section \champ{Connections}, repérez la ligne
\champ{Internal Database URL}. Cliquez sur l'icône de copie à côté.

C'est un texte très long qui commence par \saisir{postgresql://}. \textbf{Collez-le
tout de suite quelque part} — un fichier texte sur votre ordinateur — vous en
avez besoin dans deux minutes.

\begin{attention}
Cette adresse est un \textbf{mot de passe déguisé} : elle donne un accès complet
à toutes vos données. Ne l'envoyez jamais par message public, ne la publiez
nulle part, ne la mettez pas dans un email à un inconnu.
\end{attention}

% ==================================================================
\section{Étape 5 — Mettre le site en ligne}

\begin{etape}
\textbf{Dans le tableau de bord Render}
\end{etape}

\begin{enumerate}
  \item \bouton{New +}, puis \bouton{Web Service}.
  \item Render affiche vos dépôts GitHub. Choisissez
        \saisir{Fils-NAturel-Tech-com} et cliquez \bouton{Connect}.
  \item Remplissez la page de configuration :
\end{enumerate}

\begin{center}
\begin{tabularx}{\textwidth}{@{}lX@{}}
\toprule
\textbf{Champ} & \textbf{Ce qu'il faut mettre} \\
\midrule
\champ{Name} & \saisir{fntc} — ce nom apparaîtra dans l'adresse temporaire. \\
\champ{Region} & \textbf{La même que la base} (Ohio ou Virginia). C'est important. \\
\champ{Branch} & \saisir{main} \\
\champ{Runtime} & \saisir{Node} \\
\champ{Build Command} & \saisir{npm ci -{}-include=dev \&\& npm run db:init \&\& npm run build} \\
\champ{Start Command} & \saisir{npm run start} \\
\champ{Instance Type} & \textbf{Starter}, environ 7 \$/mois. Pas \champ{Free} : le site s'endormirait. \\
\bottomrule
\end{tabularx}
\end{center}

\subsection{Les deux réglages secrets}

Toujours sur la même page, section \champ{Environment Variables}, cliquez
\bouton{Add Environment Variable} \textbf{deux fois} :

\begin{center}
\begin{tabularx}{\textwidth}{@{}lX@{}}
\toprule
\textbf{Key (le nom)} & \textbf{Value (la valeur)} \\
\midrule
\saisir{DATABASE\_URL} & L'adresse copiée à l'étape 4, celle qui commence par \saisir{postgresql://} \\
\saisir{SESSION\_SECRET} & Une longue suite de caractères au hasard, au moins 32 signes \\
\bottomrule
\end{tabularx}
\end{center}

\begin{conseil}
Pour \saisir{SESSION\_SECRET}, tapez n'importe quoi de long et
d'imprévisible, par exemple :
\saisir{k7Qm2xVp9LbT4rWz8YnE3sHc6JdA5fGu}. Personne n'a besoin de s'en
souvenir : ce n'est pas un mot de passe à taper, juste une clé que le site
utilise pour signer les sessions. Ne la changez plus après coup, sinon tout le
monde est déconnecté.
\end{conseil}

\begin{enumerate}\setcounter{enumi}{3}
  \item \bouton{Create Web Service}.
  \item Render construit le site. Le journal défile pendant \textbf{cinq à dix
        minutes}. C'est normal et il n'y a rien à faire.
  \item Quand le statut passe à \champ{Live} en vert, cliquez sur l'adresse
        affichée en haut. Elle ressemble à ceci :

        \begin{center}\ttfamily https://fntc.onrender.com\end{center}

        \textbf{Le site est en ligne.}
\end{enumerate}

\subsection{Si le déploiement échoue}

\begin{center}
\begin{tabularx}{\textwidth}{@{}lX@{}}
\toprule
\textbf{Message dans le journal} & \textbf{Ce qu'il faut faire} \\
\midrule
\saisir{Environment variable not found: DATABASE\_URL} & La variable est mal écrite ou absente. Onglet \champ{Environment}, vérifiez le nom exactement : majuscules et tiret bas compris. \\
\saisir{Cannot find module} & La commande de construction est incomplète. Onglet \champ{Settings}, recopiez la \champ{Build Command} du tableau ci-dessus. \\
\saisir{Build failed} sans autre indice & Onglet \champ{Manual Deploy}, puis \bouton{Clear build cache \& deploy}. \\
\bottomrule
\end{tabularx}
\end{center}

% ==================================================================
\section{Étape 6 — Brancher le nom de domaine}

Le site fonctionne, mais à une adresse en \saisir{.onrender.com}. On y met
maintenant le vrai nom.

\subsection{Côté Render : déclarer le domaine}

\begin{enumerate}
  \item Ouvrez votre service, onglet \champ{Settings}, section
        \champ{Custom Domains}.
  \item \bouton{Add Custom Domain}. Tapez \saisir{filsnaturel.com}
        (votre domaine, sans \saisir{https://} ni \saisir{www}). Validez.
  \item Recommencez avec \saisir{www.filsnaturel.com}.
  \item Render affiche alors un tableau avec les valeurs à recopier. \textbf{Ne
        fermez pas cette page} : gardez-la ouverte dans un onglet.
\end{enumerate}

Ce tableau ressemble à ceci — \textbf{vos valeurs seront différentes,
recopiez les vôtres} :

\begin{center}
\begin{tabularx}{\textwidth}{@{}llX@{}}
\toprule
\textbf{Type} & \textbf{Nom} & \textbf{Valeur} \\
\midrule
A & \saisir{@} & \saisir{216.24.57.1} (le numéro donné par Render) \\
CNAME & \saisir{www} & \saisir{fntc.onrender.com} \\
\bottomrule
\end{tabularx}
\end{center}

\subsection{Côté Namecheap : remplir l'annuaire}

\begin{enumerate}
  \item Connectez-vous sur \lien{https://www.namecheap.com}{namecheap.com}.
  \item Menu \champ{Domain List}, puis \bouton{Manage} en face de votre domaine.
  \item Onglet \champ{Advanced DNS}.
  \item Supprimez les lignes déjà présentes qui pointent vers une page
        d'attente Namecheap : celles nommées \saisir{@} ou \saisir{www} de type
        \champ{URL Redirect} ou \champ{CNAME} vers \saisir{parkingpage}.
        Utilisez l'icône de corbeille.
  \item \bouton{Add New Record}, puis créez la première ligne :
  \begin{itemize}
    \item \champ{Type} : \saisir{A Record}
    \item \champ{Host} : \saisir{@}
    \item \champ{Value} : le numéro donné par Render
    \item \champ{TTL} : \saisir{Automatic}
  \end{itemize}
  \item \bouton{Add New Record} de nouveau, pour la deuxième ligne :
  \begin{itemize}
    \item \champ{Type} : \saisir{CNAME Record}
    \item \champ{Host} : \saisir{www}
    \item \champ{Value} : \saisir{fntc.onrender.com} (l'adresse donnée par Render)
    \item \champ{TTL} : \saisir{Automatic}
  \end{itemize}
  \item Cliquez sur la coche verte pour enregistrer chaque ligne.
\end{enumerate}

\subsection{Attendre, puis vérifier}

Retournez sur la page Render et cliquez \bouton{Verify}. Deux cas :

\begin{itemize}
  \item \textbf{Verified} en vert : c'est fait.
  \item \textbf{Pending} : l'annuaire mondial n'a pas encore été mis à jour.
        C'est normal. Attendez, puis recliquez \bouton{Verify}. Comptez dix
        minutes en général, jusqu'à 24 heures dans le pire des cas.
\end{itemize}

Une fois vérifié, Render crée tout seul le certificat HTTPS. Quelques minutes
plus tard, \saisir{https://filsnaturel.com} affiche le site avec le cadenas.
\textbf{Il n'y a rien à payer ni à renouveler pour ce cadenas.}

\begin{conseil}
L'ancienne adresse en \saisir{.onrender.com} continue de fonctionner. Elle ne
gêne pas : plus personne n'a besoin de la connaître. Communiquez uniquement le
nouveau nom.
\end{conseil}

% ==================================================================
\section{Étape 7 — Vérifier que tout marche}

Ouvrez \saisir{https://votredomaine.com} et vérifiez, dans l'ordre :

\begin{enumerate}
  \item Le cadenas apparaît dans la barre d'adresse.
  \item La page d'accueil s'affiche avec les images.
  \item Le catalogue s'ouvre, les filtres par rayon fonctionnent.
  \item Une fiche produit s'ouvre, on peut taper une quantité et l'ajouter au panier.
  \item Le panier affiche la bonne quantité et le bon total.
  \item Refaites ces cinq points \textbf{sur un téléphone}.
  \item \saisir{https://votredomaine.com/admin} affiche la page de connexion.
  \item Vous vous connectez, \textbf{et vous changez immédiatement le mot de
        passe} dans \champ{Administration $\rightarrow$ Comptes admin}.
\end{enumerate}

\begin{attention}
Le mot de passe fourni au départ est un mot de passe de livraison. Il a
circulé, ne serait-ce que dans un message. Changez-le à la première connexion,
sans exception.
\end{attention}

% ==================================================================
\section{Étape 8 — Protéger les données}

C'est l'étape que tout le monde saute, et c'est celle qu'on regrette.

\begin{enumerate}
  \item Dans Render, ouvrez la base, onglet \champ{Backups}. Vérifiez que les
        sauvegardes automatiques sont actives sur votre offre.
  \item \textbf{Une fois par mois}, faites une copie qui ne vive pas chez
        Render : bouton de téléchargement de la sauvegarde, puis rangez le
        fichier sur votre ordinateur \textbf{et} sur une clé USB ou un espace
        en ligne.
\end{enumerate}

\begin{conseil}
Une sauvegarde stockée uniquement chez l'hébergeur ne protège pas d'une erreur
de manipulation ni d'un compte fermé. Une copie chez vous, tous les mois, suffit
à dormir tranquille.
\end{conseil}

% ==================================================================
\section{Ce qu'il faut me transmettre}
\label{sec:fiche}

Remplissez cette fiche au fur et à mesure, et transmettez-la-moi. Elle me permet
de dépanner sans vous déranger. \textbf{Ne l'envoyez pas par message public} :
remettez-la en main propre, ou par un moyen privé.

\vspace{0.6em}
\begin{center}
\begin{tabularx}{\textwidth}{@{}lX@{}}
\toprule
\multicolumn{2}{@{}l}{\textbf{Le nom de domaine}}\\
\midrule
Nom exact du domaine & \aremplir \\
Acheté chez (Namecheap, Cloudflare\ldots) & \aremplir \\
Email du compte & \aremplir \\
Date d'achat & \aremplir \\
Renouvellement automatique activé ? & \aremplir \\
\midrule
\multicolumn{2}{@{}l}{\textbf{L'hébergement Render}}\\
\midrule
Email du compte Render & \aremplir \\
Adresse temporaire (\saisir{...onrender.com}) & \aremplir \\
Nom du service web & \aremplir \\
Nom de la base de données & \aremplir \\
Région choisie & \aremplir \\
\midrule
\multicolumn{2}{@{}l}{\textbf{L'administration du site}}\\
\midrule
Identifiant administrateur & \aremplir \\
Mot de passe changé le & \aremplir \\
\midrule
\multicolumn{2}{@{}l}{\textbf{Les dates à retenir}}\\
\midrule
Renouvellement du domaine (chaque année) & \aremplir \\
Jour de prélèvement Render (chaque mois) & \aremplir \\
Dernière sauvegarde téléchargée le & \aremplir \\
\bottomrule
\end{tabularx}
\end{center}

\begin{attention}
\textbf{Ne me transmettez jamais} : le mot de passe de votre compte Render, le
mot de passe de votre compte Namecheap, votre numéro de carte bancaire, ni
l'adresse \saisir{DATABASE\_URL} complète. Rien de tout cela n'est nécessaire
pour vous dépanner. Si quelqu'un vous les réclame, c'est une tentative
d'escroquerie.
\end{attention}

\begin{conseil}
Pour me donner accès sans partager de mot de passe : dans Render, section
\champ{Members} ou \champ{Team}, invitez mon adresse email. L'accès se retire
d'un clic quand il n'est plus utile.
\end{conseil}

% ==================================================================
\section{Vivre avec, mois après mois}

\begin{center}
\begin{tabularx}{\textwidth}{@{}lX@{}}
\toprule
\textbf{Quand} & \textbf{Quoi} \\
\midrule
Chaque mois & Render prélève environ 13 \$. Vérifiez que le paiement est passé : une carte expirée coupe le site. \\
Chaque mois & Téléchargez une sauvegarde de la base et rangez-la chez vous. \\
Chaque année & Le domaine se renouvelle. Si votre carte a changé, mettez-la à jour chez le registrar \textbf{avant} l'échéance. \\
Quand le site évolue & Un \saisir{git push} sur GitHub déclenche un nouveau déploiement tout seul. Vous n'avez rien à faire. \\
\bottomrule
\end{tabularx}
\end{center}

\begin{attention}
\textbf{Le danger le plus courant n'est pas technique : c'est une carte
expirée.} Un domaine non renouvelé peut être racheté par n'importe qui, et se
récupérer coûte alors très cher — quand c'est encore possible. Surveillez les
emails de Namecheap et de Render, ils préviennent à l'avance.
\end{attention}

% ==================================================================
\section{Si quelque chose ne va pas}

\begin{center}
\begin{tabularx}{\textwidth}{@{}lX@{}}
\toprule
\textbf{Symptôme} & \textbf{Que faire} \\
\midrule
Le domaine affiche « site introuvable » & L'annuaire n'est pas encore à jour. Attendez, jusqu'à 24 h. Vérifiez ensuite que les deux lignes DNS sont exactement celles données par Render. \\
Le domaine affiche une page de publicité Namecheap & Il reste une ancienne ligne \champ{URL Redirect} ou \champ{parkingpage} dans \champ{Advanced DNS}. Supprimez-la. \\
Pas de cadenas, « connexion non sécurisée » & Le certificat n'est pas encore émis. Attendez quinze minutes après la vérification du domaine. \\
Le site met vingt secondes à s'ouvrir & Vous êtes resté sur l'offre gratuite. Passez le service en \champ{Starter}. \\
Le site affiche une erreur serveur & Onglet \champ{Logs} du service : la dernière ligne rouge dit pourquoi. Envoyez-la-moi telle quelle. \\
Le site a disparu, la base est vide & L'offre gratuite de base de données a expiré. C'est irréversible sans sauvegarde. \\
Plus rien ne répond, aucune erreur & Vérifiez le paiement : une carte refusée suspend le service. \\
\bottomrule
\end{tabularx}
\end{center}

\subsection{Où demander de l'aide}

\begin{itemize}
  \item Render : \lien{https://render.com/docs}{render.com/docs} et le chat de support depuis le tableau de bord.
  \item Namecheap : chat en direct 24\,h/24 depuis le site, en anglais.
  \item Pour tout ce qui touche au site lui-même : passez par moi, avec la
        ligne d'erreur exacte et une capture d'écran.
\end{itemize}

\vspace{1.2cm}
\begin{center}
{\color{sig}\rule{3cm}{1.5pt}}\par
\vspace{0.4cm}
{\color{encre2}\small Fils Naturel Tech-Com — guide de mise en ligne}
\end{center}

\end{document}
